% authority: science_draft
\documentclass[11pt]{article}

\usepackage[margin=1in]{geometry}
\usepackage{amsmath,amssymb}
\usepackage{booktabs}
\usepackage{longtable}
\usepackage{enumitem}

\newcommand{\EqSrc}{\mathrm{EqSrc}}
\newcommand{\EqSrcT}{\mathrm{EqSrc}_{T}}
\newcommand{\RetainH}{\mathrm{RetainH}}
\newcommand{\GenH}{\mathrm{GenH}}
\newcommand{\Fsrc}{\mathcal{F}_{\mathrm{src}}}
\newcommand{\Osrc}{\mathcal{O}_{\mathrm{src}}}
\newcommand{\Msrc}{\mathcal{M}_{\mathrm{src}}}
\newcommand{\Isrc}{\mathcal{I}_{\mathrm{src}}}
\newcommand{\Asrc}{\mathcal{A}_{\mathrm{src}}}
\newcommand{\Psrc}{\mathcal{P}_{\mathrm{src}}}
\newcommand{\Nsrc}{\mathcal{N}_{\mathrm{src}}}
\newcommand{\Csrc}{C_{\mathrm{src}}}
\newcommand{\NoTarget}{\mathsf{NoTarget}}

\title{EqSrc Family-Closure Smuggling Audit v1}
\author{Smuggling Auditor Packet RT-20260707-022}
\date{2026-07-07}

\begin{document}
\maketitle

\section{Control Status}

This artifact implements v18 P3-T04 as a draft/control smuggling audit. It
audits the P3-T02 conditional \(\EqSrcT\) family-closure theorem candidate and
its missing-inverse countermodel slot for target imports and authority
laundering.

The audit result is:
\[
  \texttt{source\_pure\_as\_written}.
\]

\begin{verbatim}
audit_result: source_pure_as_written
family_closure_status: supplied_as_conditional_H1_H7_hypotheses_not_derived
countermodel_audit: source_side_as_written
next_route: P3-T05
\end{verbatim}

This result has no proof-promotion authority. It does not discharge general
\(\EqSrc\), adopt \(\RetainH\), adopt \(\GenH\), adopt a source law, edit
canonical ontology, adopt \(\mathrm{MetricData}(E)\), expand
\(g_{\mathrm{eff}}\), derive matter coupling, derive Einstein equations,
promote a benchmark, issue a Gate Chair verdict, authorize external outreach,
or claim completed derivation.

\section{Artifact Under Audit}

The audited artifact is:
\begin{quote}
\texttt{research\_control/tasks/RT-20260707-020/artifacts/eqsrc\_family\_closure\_theorem\_or\_countermodel\_v1.tex}.
\end{quote}

The audit also inspects the primitive-boundary extraction:
\begin{quote}
\texttt{research\_control/tasks/RT-20260707-021/artifacts/retainh\_genh\_primitive\_boundary\_v1.tex}.
\end{quote}

The P3-T02 primary branch is
\texttt{family\_closure\_theorem\_candidate\_supplied}. It states that a
declared source family \(\Fsrc\) satisfies hypotheses H1--H7 and then proves
reflexivity, symmetry, and transitivity for \(\EqSrcT\). The P3-T03 artifact
records that \(\RetainH\) and \(\GenH\) are not required for the already
declared closed-family theorem and remain candidate-definition-needed for
\(H\)-retention or \(H\)-generated extensions.

\section{Audit Method}

For each possible smuggling channel, the audit asks whether the theorem branch
or countermodel branch uses that channel as a premise. Mentions inside
explicit forbidden-premise guards count as guards, not as imports.

The audit distinguishes three statuses:
\begin{itemize}[leftmargin=*]
  \item \texttt{source\_pure\_as\_written}: no target or process-authority
  import is detected in the written artifact.
  \item \texttt{repair\_required}: a premise is ambiguous enough that the
  artifact must be rewritten before stress.
  \item \texttt{fails\_closed}: the artifact depends on a forbidden target or
  process-authority premise.
\end{itemize}

\section{Theorem Branch Audit: Supplied versus Derived}

The theorem branch is source-pure as written because its proof only uses the
following declared source-side hypotheses:
\begin{enumerate}[label=(H\arabic*)]
  \item source-only object domain;
  \item identity closure;
  \item inverse closure;
  \item composition closure;
  \item invariant-ledger stability;
  \item source-only comparison;
  \item no-target guard.
\end{enumerate}

The key audit finding is not that family closure has been derived from current
ontology. The key finding is that family closure was supplied conditionally:

\begin{verbatim}
family_closure_derived_from_current_ontology: false
family_closure_supplied_as_conditional_hypotheses: true
H1_H7_status: supplied_not_derived
\end{verbatim}

Therefore the theorem branch is source-pure as a conditional theorem
candidate, but it cannot be read as a general \(\EqSrc\) discharge. A later
packet must still stress H1--H7, especially inverse closure, composition
closure, ledger stability, and no-target robustness.

\section{Countermodel Branch Audit}

The required countermodel slot in P3-T02 is the missing-inverse countermodel.
It uses a finite source object set \(\{A,B\}\), a source morphism
\(f:A\to B\), identity morphisms, and a source comparison rule that blocks
\(\Csrc(B,A)\) because no source-side inverse witness is available.

The countermodel branch is:

\begin{verbatim}
countermodel_branch_checked: true
countermodel_branch_finding: source_side_as_written
global_no_go_claimed: false
\end{verbatim}

It does not import target topology, target metric, detector protocol,
empirical calibration, stress-energy, matter action, Einstein equations,
benchmark behavior, or process authority. It also does not prove global
impossibility. It only records that inverse closure is a real hypothesis.

\section{Target-Import Audit Matrix}

\begin{longtable}{p{0.27\linewidth}p{0.17\linewidth}p{0.48\linewidth}}
\toprule
Audit target & Finding & Reason \\
\midrule
target topology import & not detected & The theorem and countermodel use source objects, source morphisms, and source comparison data only. \\
target smooth atlas import & not detected & No target charts, transition maps, differentiability, or atlas compatibility condition is a premise. \\
target Lorentzian metric import & not detected & No metric tensor, signature, connection, curvature, cone, or geodesic structure supplies the proof. \\
proper-time normalization import & not detected & No proper-time clock or worldline normalization enters H1--H7. \\
detector protocol import & not detected & Detector protocol language is absent as a premise and remains a forbidden downstream channel. \\
empirical calibration import & not detected & No measurement table, calibration law, or empirical response supplies source equivalence. \\
stress-energy semantics import & not detected & No stress-energy tensor, density, pressure, flux, or energy condition is used. \\
matter-action import & not detected & No matter action or Lagrangian density is used. \\
Einstein-equation premise import & not detected & No Einstein tensor, field equation, or field-equation residual is used. \\
benchmark behavior import & not detected & Exact-GR benchmark behavior is not a theorem premise. \\
\bottomrule
\end{longtable}

\section{Process-Authority Audit Matrix}

\begin{longtable}{p{0.29\linewidth}p{0.17\linewidth}p{0.46\linewidth}}
\toprule
Audit target & Finding & Reason \\
\midrule
validation status as premise & not detected & Validator passes are operational receipts only and do not establish H1--H7. \\
registry metadata as theorem premise & not detected & Registry rows route and identify artifacts; they do not make source morphisms exist. \\
role identity as theorem premise & not detected & The authoring role constrains execution but is not a mathematical input. \\
handoff state as theorem premise & not detected & Handoffs select packets and preserve routing; they do not prove family closure. \\
commit state as theorem premise & not detected & Commits preserve provenance only. \\
generated derivative as theorem premise & not detected & Wiki notes, indexes, semantic extracts, and local memory remain derivative retrieval layers. \\
\bottomrule
\end{longtable}

\section{Family-Closure Source-Purity Lemma}

\textbf{Lemma.} If the P3-T02 theorem candidate is read exactly as written,
with H1--H7 treated as source-side conditional hypotheses and with
\(\NoTarget\) active, then target topology, target smooth atlas structure,
target Lorentzian metric structure, proper-time normalization, detector
protocol, empirical calibration, stress-energy semantics, matter-action
semantics, Einstein-equation premises, benchmark behavior, validator status,
registry metadata, role identity, handoff state, commit state, and generated
derivatives are not premises of the \(\EqSrcT\) family-closure proof.

\textbf{Reasoning.} Reflexivity uses H1, H2, H5, H6, and H7. Symmetry uses H3,
H5, H6, and H7. Transitivity uses H4, H5, H6, and H7. Each hypothesis is
stated as source-side in the artifact. The proof does not appeal to a target
metric, target atlas, detector protocol, empirical calibration, field
equation, benchmark result, validator, registry row, role, handoff, commit, or
generated derivative. Therefore the theorem branch passes the smuggling audit
as a conditional source theorem candidate.

\section{Audit Result}

The audit result is:

\begin{center}
\texttt{source\_pure\_as\_written}
\end{center}

The exact scope is a hidden-import screen for the P3-T02 conditional theorem
candidate and missing-inverse countermodel slot. The allowed use is routing to
P3-T05 Refuter stress. The blocked overread is any claim that current ontology
derives H1--H7, any general \(\EqSrc\) discharge, and any RetainH or GenH
adoption.

\section{Residual Risks and Repairs}

No repair is required before P3-T05. The stress packet should nevertheless
target the strongest remaining pressure points:
\begin{itemize}[leftmargin=*]
  \item inverse closure can fail when a source-side inverse witness is absent;
  \item composition closure can fail when composites leave \(\Msrc\);
  \item ledger stability can fail under family variation;
  \item \(\RetainH\) remains only a candidate-definition-needed pressure for
  \(H\)-retention extensions;
  \item \(\GenH\) remains only a candidate-definition-needed pressure for
  \(H\)-generated family extensions;
  \item a local audit pass can be overread as proof authority unless P3-T05
  stress preserves the supplied-not-derived boundary.
\end{itemize}

\section{Distance-to-GR Effect}

\begin{longtable}{p{0.30\linewidth}p{0.56\linewidth}}
\toprule
Burden & Status after this packet \\
\midrule
Source ontology primitives & Unchanged; no canonical ontology edit. \\
Source equivalence EqSrc & Conditional family-closure theorem candidate audited source-pure as written; no general EqSrc discharge. \\
RetainH & Not adopted; still not required for the closed-family theorem and candidate-definition-needed for H-retention extensions. \\
GenH & Not adopted; still not required for the closed-family theorem and candidate-definition-needed for H-generated extensions. \\
ObsLoc\_lc & Unchanged. \\
Resp\_lc & Unchanged. \\
M\_src & Unchanged. \\
g\_eff & Unchanged; no MetricData(E) or metric-scope expansion. \\
matter coupling & Blocked and unchanged; no coupling derivation or adoption. \\
Einstein equations & Blocked and unchanged; no field-equation derivation. \\
finite-variation robustness & Unchanged; P3-T05 stress remains required. \\
benchmark promotion & Blocked; no benchmark evidence or promotion. \\
Gate Chair status & Unchanged; no Gate Chair verdict. \\
current route freeze or hard-fail status & Not frozen; route continues to P3-T05 stress. \\
\bottomrule
\end{longtable}

\begin{verbatim}
distance_to_gr_delta:
  changed: false
  effect: no_distance_delta
  ledger_row_updated: false
  reason: P3-T04 is a hidden-import audit only. It does not prove H1-H7 from
    current ontology and does not adopt EqSrc RetainH GenH or any downstream
    physics claim.
\end{verbatim}

\section{Forbidden Conclusions}

This audit result must not be read as any of the following:
\begin{itemize}[leftmargin=*]
  \item canonical ontology edit;
  \item general \(\EqSrc\) discharge;
  \item \(\RetainH\) adoption;
  \item \(\GenH\) adoption;
  \item source-law adoption;
  \item source detector/readout semantics adoption;
  \item \(\mathrm{MetricData}(E)\) adoption;
  \item \(g_{\mathrm{eff}}\) adoption or scope expansion;
  \item matter-coupling derivation or adoption;
  \item stress-energy semantics, stress-energy tensor, or matter action;
  \item Einstein-equation derivation;
  \item benchmark promotion;
  \item Gate Chair verdict;
  \item external outreach;
  \item completed derivation;
  \item future source-extension impossibility;
  \item global no-go conclusion.
\end{itemize}

\section{Next Route}

The logical next route is one bounded v18 P3-T05 Refuter stress packet against
the audited family-closure attempt. The stress packet should target inverse
closure removal, composition closure removal, ledger-instability witnesses,
RetainH and GenH primitive-dependence pressure, no-target-guard weakening, and
overread from local source-purity to theorem or adoption status.

\section{Source Materials}

The AEther-Flow Research Project. (2026a). \emph{Recommendations
implementation plan continue task v18} [Internal implementation plan].
\texttt{implementations\_plans/recommendations\_implementation\_plan\_continue\_task-v18.md}.

The AEther-Flow Research Project. (2026b). \emph{EqSrc family-closure theorem
or countermodel v1} [Research-control TeX artifact].
\texttt{research\_control/tasks/RT-20260707-020/artifacts/eqsrc\_family\_closure\_theorem\_or\_countermodel\_v1.tex}.

The AEther-Flow Research Project. (2026c). \emph{RetainH and GenH
primitive-boundary extraction v1} [Research-control TeX artifact].
\texttt{research\_control/tasks/RT-20260707-021/artifacts/retainh\_genh\_primitive\_boundary\_v1.tex}.

The AEther-Flow Research Project. (2026d). \emph{Handoff 0690} [Internal
research-control handoff].
\texttt{research\_control/handoffs/handoff-0690.yaml}.

\end{document}
